\begin{verbatim}
a) Euler’s Method:
#include <stdio.h>

double f(double x, double y) { return x+y; }

int main() {
    double x0, y0, h, xn, x, y;
    printf("Enter x0, y0, h, xn: ");
    scanf("%lf %lf %lf %lf", &x0, &y0, &h, &xn);

    x = x0; y = y0;
    while(x < xn) {
        y = y + h*f(x,y);
        x += h;
    }
    printf("Approximate y(%lf) = %lf\n", xn, y);
    return 0;
}

b) Runge-Kutta 2nd & 4th Order:
#include <stdio.h>

double f(double x, double y) { return x+y; }

int main() {
    double x0, y0, h, xn, x, y, k1, k2, k3, k4;
    printf("Enter x0, y0, h, xn: ");
    scanf("%lf %lf %lf %lf", &x0, &y0, &h, &xn);

    x = x0; y = y0;
    while(x < xn) {
        // RK4
        k1 = h*f(x,y);
        k2 = h*f(x+h/2, y+k1/2);
        k3 = h*f(x+h/2, y+k2/2);
        k4 = h*f(x+h, y+k3);
        y = y + (k1+2*k2+2*k3+k4)/6;
        x += h;
    }
    printf("RK4 Approximate y(%lf) = %lf\n", xn, y);
    return 0;
}
\end{verbatim}